\documentclass[11pt]{letter}
\usepackage[T1]{fontenc}
\usepackage[a4paper,margin=0.9in]{geometry}
\usepackage{hyperref}
\signature{Hikaru Wakaura\\ \emph{on behalf of both authors}}
\address{QIRI, Quantum Integrated Research Institute Inc.\\ Tokyo 107-0061, Japan\\ \texttt{h.wakaura@qiri.co.jp}}
\begin{document}
\begin{letter}{The Editors\\ \emph{Physical Chemistry Chemical Physics}}
\opening{Dear Editors,}

On behalf of my co-author, I am pleased to submit our manuscript,
\textbf{``Ultrafast recombination and exchange coupling preclude weak-field
magnetosensitivity of neurotransmitter-metabolising flavoenzymes,''} for
consideration in \textit{Physical Chemistry Chemical Physics}.

The manuscript sits squarely in the spin-chemistry scope of PCCP. Using an exact
Liouville-space radical-pair master equation that includes spin-selective
recombination, exchange and dipolar coupling, hyperfine interactions, and
electronic decoherence, we establish a quantitative bound on whether a weak,
static magnetic field can bias the catalytic flux of the neurotransmitter-metabolising
flavoenzymes (monoamine oxidase A and B, D-amino-acid oxidase) through the
radical-pair mechanism. Two independent active-site properties---a gigahertz
exchange coupling and a picosecond recombination lifetime---each suppress the
geomagnetic-field effect below the numerical resolution of the solver, and we show
the null is robust to the form of the electronic decoherence (including the
singlet--triplet dephasing channel appropriate to a tightly bound pair) and to the
quantum Zeno mechanism. Cryptochrome, treated identically, reproduces the expected
percent-level sensitivity as a positive control. The result quantifies a
methodological boundary of broad relevance to the community: an RPM estimate for an
active-site cofactor radical must propagate the physical lifetime and electronic
decoherence, or it overestimates the weak-field effect by orders of magnitude. The
prediction is consistent with the one experiment to scan static fields on brain
monoamine oxidase, which found no change from 0.1 to 340~mT.

This is a complete study, developed independently and written from scratch around
the master-equation formalism; it is not under consideration elsewhere and has no
prior publication. All code is openly released and the central claims are covered
by an automated test suite. The authors declare no competing interests and the
study received no specific funding. We suggest the manuscript is well matched to
referees working in radical-pair spin dynamics and quantum biology.

Thank you for considering our submission.

\closing{Yours sincerely,}
\end{letter}
\end{document}
